% Source: https://artofproblemsolving.com/community/c3h2117093_drawing_fries
% Stable identity: aops-topic-2117093-q01
Imagine that the drawing continues until all $30$ fries have been drawn; this does not change who draws the longest fry. The order in which the fries come out of the box is a uniformly random ordering of the $30$ fries, so the longest fry is equally likely to be the $1$st, $2$nd, $\dots$, $30$th fry drawn.

The first $15$ fries are drawn by the pals. After that the turns go in rounds of six, starting with Jefferson, so Jefferson draws fries number
\[16,\qquad 22,\qquad 28,\]
and the pals draw all the others. Jefferson gets the coke exactly when the longest fry is in one of these three positions, so the probability is
\[\dfrac{3}{30}=\boxed{\dfrac{1}{10}}.\]
(If only Jefferson's first fry counted, the probability would be $\dfrac{1}{30}$; given that none of the first $15$ fries was the longest, his first fry would be the longest with probability $\dfrac{1}{15}$.)
